% ECONOMICS_PROBLEM: {"id": "C72-142", "title": "Strong approximation of Nash equilibria in concise potential games", "jel_code": "C72", "source_id": "OP-0140"}
% !TeX program = xelatex
\documentclass[11pt,a4paper]{article}
\usepackage[margin=23mm]{geometry}
\usepackage{fontspec}
\setmainfont{Latin Modern Roman}
\usepackage{amsmath,amssymb,mathtools}
\usepackage{xcolor,enumitem,booktabs,longtable,array,xurl,hyperref}
\definecolor{muted}{HTML}{53636C}
\hypersetup{colorlinks=true,urlcolor=blue,linkcolor=blue}
\setlength{\parindent}{0pt}
\setlength{\parskip}{5pt}
\newcommand{\R}{\mathbb R}
\newcommand{\E}{\mathbb E}
\newcommand{\N}{\mathbb N}
\newcommand{\ind}{\mathbf 1}
\newcommand{\Var}{\operatorname{Var}}
\newcommand{\Cov}{\operatorname{Cov}}
\newcommand{\supp}{\operatorname{supp}}
\DeclareMathOperator*{\argmin}{arg\,min}
\DeclareMathOperator*{\argmax}{arg\,max}
\newcommand{\entrymeta}[3]{{\sffamily\small\color{muted}\textbf{\texttt{#1}}\quad|\quad #2\par #3\par}}
\newcommand{\Problemref}[1]{\texttt{#1}}
\newcommand{\JELref}[2]{\texttt{#2} (JEL \texttt{#1})}
\begin{document}
\section*{Strong approximation of Nash equilibria in concise potential games}
\textbf{Problem ID:} \texttt{C72-142}\quad \textbf{JEL:} \texttt{C72}\par
% BEGIN ECONOMICS STATEMENT
\label{problem:OP-0140}
% GAME_THEORY_PROBLEM: {"local_id": "GT-010", "original_number": 10, "kind": "P", "entry_kind": "statement_only", "published": false, "review_required": true, "category": "game_theory"}
\label{game:GT-010}
{\sffamily\small\color{muted}
\textbf{GT-010} | Potential games and approximation complexity\\
\textbf{P} | Second explicit question in the same source.
\par}
\subsection*{Definitions and mathematical statement}
Use explicitly listed action sets and a rational multilinear potential circuit with the encoding conventions of GT-009. Let $X=\prod_i\Delta(S_i)$ and let $E\subseteq X$ be its set of exact Nash equilibria, defined by $\Phi(z)\geq\Phi(p_i,z_{-i})$ for all unilateral mixed deviations. For rational input $\varepsilon\in(0,1)$, the strong-approximation search problem asks for a rational vector $x\in X$ such that
\[
 \exists z\in E,\qquad \|x-z\|_\infty\leq\varepsilon.
\]
The precision parameter is encoded in binary, and its bit length is part of the input. Under the reductions and PLS convention specified in [S1], is this problem $\mathsf{FIXP}_a\cap\mathsf{PLS}$-complete? Here $\mathsf{FIXP}_a$ is the class of strong finite-precision approximations to exact algebraic fixed-point search problems; proximity is to an actual fixed point, not merely a small fixed-point residual.
\subsection*{Short English statement}
Classify the difficulty of outputting a point close to a genuine equilibrium in a compactly represented potential game.
\subsection*{Variants and scope boundaries}
The target permits closeness to pure equilibria. A small unilateral-payoff residual does not provide a uniform geometric-distance guarantee without a quantitative game-dependent argument. Fully expanded normal-form potential games are a different input class. Changing the norm is manageable in a polynomial-dimensional space only after the corresponding accuracy conversion is stated.
\subsection*{Source relationship and status}
The second sentence of Question 6.1 poses this strong-approximation completeness target. The paper's express caveat about PLS and verification applies here as well. No claim is made that strong approximation has efficiently checkable certificates under conventional exact-equilibrium verification.
\subsection*{References}
\begingroup\footnotesize\raggedright
{\footnotesize\raggedright
\textbf{[S1]} Ioannis Anagnostides, Maria-Florina Balcan, Kiriaki Fragkia, Tuomas Sandholm, Emanuel Tewolde, and Brian Hu Zhang (2026 version). \emph{The Complexity of Equilibrium Refinements in Potential Games}. arXiv:2511.03968. Locator: Section 3.1; Question 6.1 and footnote 5, printed p.\ 20. \url{https://arxiv.org/pdf/2511.03968}\par}
\par\endgroup

\subsection*{Common conventions: Source mathematical conventions}

\textbf{Finite games.} For a finite set $A$, $\Delta(A)$ is its probability
simplex. Players choose independent mixed strategies in Nash equilibrium;
correlation is allowed only when explicitly part of the solution concept.
In a game with payoff functions $u_i$ and strategy profile $x$, an additive
$\varepsilon$ Nash equilibrium satisfies
\[
 u_i(x)\geq u_i(x_i',x_{-i})-\varepsilon
 \quad\text{for every player $i$ and every admissible deviation $x_i'$}.
\]
For numerical approximation, payoffs are normalized as locally specified.
An $\varepsilon$ payoff residual is distinct from distance at most
$\varepsilon$ to the set of exact equilibria. Point convergence, convergence
of empirical play, and convergence in distance to an equilibrium set are
also distinct.

\textbf{Computational representations.} Unless stated otherwise, a complexity
question takes rational data with binary encoding; polynomial time is measured
in total bit length. A fixed-accuracy polynomial algorithm is not necessarily
polynomial in $1/\varepsilon$, and neither is necessarily polynomial in
$\log(1/\varepsilon)$. Succinct games and value, demand, or sampling oracles
must declare their interfaces. Exact real solutions require an explicit
representation; an unspecified list of arbitrary real numbers is not a
finite computational output. A claimed algorithmic barrier is meaningful
only for its specified model and reductions.

\textbf{Dynamic games.} Histories, information, observation, and allowable
randomization are part of the model. A stationary strategy depends only on
the current state; a positional strategy in a deterministic graph game
chooses one action at each controlled vertex. Nash equilibrium at an initial
state and equilibrium after every history are different requirements.
An expectation of a pathwise $\liminf$ payoff, a $\liminf$ of expected averages,
a limiting expected average, and a uniform finite-horizon equilibrium are
different criteria. None is silently substituted for another.

\textbf{Allocation and mechanisms.} A complete allocation assigns every item
exactly once unless otherwise stated. Zero-valued goods, negative values,
capacity constraints, feasibility, lotteries, and copies of goods affect
fairness definitions and are specified locally. Dominant-strategy incentive
compatibility, Bayesian incentive compatibility, universal truthfulness,
and truthfulness in expectation impose different inequalities. Whenever
an objective ratio has zero denominator, its convention is stated or those
instances are excluded.

\textbf{Infinite-dimensional models.} Expectations, integrals, stochastic
equations, and optimization objectives require their local regularity,
integrability, filtrations, and admissible domains. A backward--forward
mean-field system is a boundary-value problem; it is not an initial-value
semiflow without an additional construction. Long-time, large-population,
and vanishing-noise limits require an order of limits and a convergence
topology. If a minimum need not be attained, the objective is its infimum
or supremum and the approximation requirement is stated separately.
% END ECONOMICS STATEMENT
\end{document}
